\section{Runtime and memory accounting}
\label{app:complexity-details}

The runtime of a fit is the sum, over realized nodes, of the scans, copies, and
tests each node reaches. Feature pools, node sizes, candidate counts, and
stopping times all vary across the tree, so the accounting is node by node.

\paragraph{Notation.}
Let $n$ be the training sample count and $p$ the number of features. At node
$t$, $n_t$ samples arrive, $a_t=|F_{t,\mathrm{avail}}|$ features are available,
$d_t^{\mathrm{const}}$ of them are constant at the node, and Stage~A considers
$m_t=|F_t|$ features after pruning and subsampling. Feature muting removes
$d_t^{\mathrm{mute}}$ tested features from the descendant pool. When Stage~B is
reached, $\tilde{K}_t$ thresholds are generated for $j_t^\star$ and
$K_t\le\tilde{K}_t$ remain after the minimum leaf size filter. Let
$C^{\mathrm{sel}}_j(n_t)$ be the cost of one Stage~A statistic and
$C^{\mathrm{split}}(n_t)$ the cost of one weighted child impurity, excluding
label copies, shuffles, and masks. Both are linear in $n_t$ for the MC selector
with Gini splitting and the PC selector with MSE splitting. The RDC costs
$O(n_t\log n_t)$ for fixed projection count, and distance correlation costs
$O(n_t\log n_t)$ or $O(n_t^2)$ depending on the backend.

\paragraph{Where each runtime control enters.}
Feature subsampling caps $m_t$, the Stage~A testing width, but not the
constant-feature pass over all $a_t$ available features. Feature muting shrinks
$a_u$ and $m_u$ at descendants $u$ at a cost of $O(a_t d_t^{\mathrm{mute}})$
index updates at $t$. Histogram and other bounded threshold methods reduce
$\tilde{K}_t$ when their cap is below the number of midpoints; exact search and
features with at most four unique values use all midpoints. Adaptive stopping,
feature scanning, and threshold scanning reduce work only when they shorten the
realized path; the worst case is still every feature, every threshold, and the
full budget, and scanning pays for a prescan. With the corrections enabled, a
stage with $n_{\mathrm{tests}}$ tests ($m_t$ in Stage~A, $K_t$ in Stage~B) at
nominal level $\alpha_{\mathrm{stage}}$ tests each at
$\alpha_{\mathrm{adj}}=\alpha_{\mathrm{stage}}/n_{\mathrm{tests}}$. The minimum
budget is then $\lceil n_{\mathrm{tests}}/\alpha_{\mathrm{stage}}\rceil$, the
maximum budget is about $n_{\mathrm{tests}}^2/(4\alpha_{\mathrm{stage}}^2)$,
and the auto budget of \cref{eq:auto-budget} is 52 at
$\alpha_{\mathrm{adj}}=0.05$, 3{,}144 at $0.05/20$, and 64{,}669 at $0.05/257$.
An explicit integer budget is multiplied by $n_{\mathrm{tests}}$. Forest size
sets the number of tree fits, and the bootstrap sample size sets the root size
of each tree.

\paragraph{Per-node work.}
Every node pays $O(n_t)$ for the stopping checks and leaf value. A node that can
split scans the available features for constants, at $O(a_t n_t)$, and removes
them, at $O(a_t d_t^{\mathrm{const}})$. Let $\mathcal{J}_t$ be the Stage~A
features whose test is evaluated before the scan stops, and
$b^{\mathrm{sel}}_{t,j}$ the permutations drawn for feature $j$. The Stage~A
work is
\begin{equation}
  S^{\mathrm{A}}_t
  = O\!\left(\sum_{j \in \mathcal{J}_t}
      \bigl(b^{\mathrm{sel}}_{t,j}+1\bigr)
      \bigl(C^{\mathrm{sel}}_j(n_t)+n_t\bigr)\right),
\end{equation}
where the $+1$ is the observed statistic and the $n_t$ term covers label
copies and shuffles. Feature scanning adds
$O(\sum_{j\in F_t} C^{\mathrm{sel}}_j(n_t) + m_t\log m_t)$ for the prescan and
sort.

Stage~B builds the thresholds of the selected feature in $O(n_t\log n_t)$,
because unique values are sorted before any cap applies. Let
$\mathcal{C}^{\mathrm{eval}}_t$ be the thresholds tested before the scan stops,
and $b^{\mathrm{split}}_{t,c}$ the permutations drawn at threshold $c$. The
Stage~B work is
\begin{equation}
  S^{\mathrm{B}}_t
  = O\!\left(\sum_{c \in \mathcal{C}^{\mathrm{eval}}_t}
      \bigl(b^{\mathrm{split}}_{t,c}+1\bigr)
      \bigl(C^{\mathrm{split}}(n_t)+n_t\bigr)\right),
\end{equation}
where the $n_t$ term covers masks, child-label slices, copies, and shuffles.
The two Stage~B rules differ in how the budget enters this sum. The
per-threshold rule tests each threshold at $\alpha_{\mathrm{split}}/K_t$, so at
the minimum budget each test draws $\lceil K_t/\alpha_{\mathrm{split}}\rceil$
permutations and
\[
  S^{\mathrm{B}}_t
  = O\!\left(|\mathcal{C}^{\mathrm{eval}}_t|\,
      \lceil K_t/\alpha_{\mathrm{split}}\rceil
      \bigl(C^{\mathrm{split}}(n_t)+n_t\bigr)\right),
\]
which is quadratic in the candidate count when every threshold is tested. The
max-type rule draws one set of $B$ permutations for all $K_t$ candidates, so
\[
  S^{\mathrm{B}}_t
  = O\!\left((B+1)\,K_t\,\bigl(C^{\mathrm{split}}(n_t)+n_t\bigr)\right),
\]
which is linear in the candidate count at fixed $B$. Threshold scanning adds
$O(K_t(C^{\mathrm{split}}(n_t)+n_t) + K_t\log K_t)$ for ordering. An accepted
split copies the node data into two child matrices with all $p$ columns, at
$O(n_t p)$. Nodes that become leaves after Stage~A or Stage~B keep the work
already done and allocate no children.

\paragraph{Trees and forests.}
Let $W_t$ be the work at node $t$. A tree costs
$T_{\mathrm{tree}}=\sum_{t\in\mathcal{V}(\mathcal{T})} W_t$ over all realized
nodes, leaves included, and a binary tree has
$|\mathcal{V}(\mathcal{T})|=2|\mathcal{I}(\mathcal{T})|+1$ nodes for
$|\mathcal{I}(\mathcal{T})|$ internal splits. Honesty computes the node sum on
the structure subsample and adds $O(np)$ for the partition plus a leaf
re-estimation pass. For a forest of $M$ trees, with $s_r$ samples for tree $r$
after bootstrap sampling,
\begin{equation}
  T_{\mathrm{forest}}
  = O(np) + \sum_{r=1}^{M}\bigl(O(s_r p) + T_{\mathrm{tree},r}\bigr) + O(Mp),
\end{equation}
plus the honesty and out-of-bag scoring passes when they are requested. This
is total work; wall-clock time also depends on the worker count and load
balance across trees.

\paragraph{Memory.}
Node-local workspaces are $O(a_t+m_t)$ for feature indices and
$O(n_t+\tilde{K}_t)$ for thresholds, plus the permutation buffers; a
fixed-budget test without stopping may store its $B$ permuted statistics. The
dense matrices dominate. Validation materializes an $O(np)$ matrix, and each
accepted split materializes children of combined size $O(n_t p)$. During
depth-first recursion the parent keeps the sibling child alive, so with
$\mathcal{P}$ the accepted splits on the active recursion stack, peak
single-tree memory is
\[
  O\!\left(np + p\sum_{u\in\mathcal{P}} n_u\right),
\]
which is $O(np)$ for balanced growth. A forest runs up to
$w=\min(M,n_{\mathrm{jobs}})$ trees at once, so its peak memory holds $w$ such
workspaces. Fitted storage holds the node records of all trees and $O(Mp)$
per-tree arrays.

\subsection{Measured scaling in \texorpdfstring{$n$}{n} and \texorpdfstring{$p$}{p}}
\label{app:scaling-curves}

\Cref{tab:scaling-curves} reports fit-only time for one CIT and one 100-tree CIF
on Gaussian predictors with a weak planted signal. The configuration uses the
MC selector with Gini splitting in classification and the PC selector with MSE
splitting in regression, $\alpha_{\mathrm{sel}}=\alpha_{\mathrm{split}}=0.05$
with both corrections, the minimum budget, the per-threshold rule, the 256-bin
histogram, feature and threshold scanning, and feature muting. The forest adds
bootstrap sampling and square-root feature sampling.

The histogram caps the candidates at 257 and the minimum budget fixes at most
5{,}140 permutations per threshold, so the per-node cost is linear in $n_t$. A
64-fold increase in $n$ at $p=20$ costs the tree about 68 times in
classification and 110 times in regression. A 25-fold increase in $p$ at
$n=2{,}000$ costs it 21 and 44 times, because every predictor is scored at the
root and muting removes only failed predictors. The forest grows 65 to 77 times
for the 64-fold increase in $n$ and is flat or falling in $p$ under square-root
feature sampling.

Adaptive stopping is inert at the minimum budget (\cref{prop:adaptive-size}),
but its checks still add time. The adaptive-to-no-stopping time ratio is 1.02 to 1.28 for the
tree (median 1.07) and 0.95 to 1.18 for the forest (median 0.99). The adaptive
path runs through the chunked sequential kernel, which carries a fixed per-test
overhead. The tree overhead grows with $p$, from 1.02 to 1.07 at $p=20$ to 1.15
to 1.28 at $p=500$, because a single tree runs one Stage~A test per predictor
at every node. These ratios match \cref{tab:runtime-ablations}.

\begin{table}[!htbp]
  \centering
  \footnotesize
  \setlength{\tabcolsep}{5pt}
  \begin{tabular}{@{}lrrrrrr@{}}
    \toprule
     & & & \multicolumn{2}{c}{CIT} & \multicolumn{2}{c}{CIF} \\
    \tablehead{Task} & $n$ & $p$ & Adaptive & No stopping & Adaptive & No stopping \\
    \midrule
    Classification & 500 & 20 & 0.17 & 0.16 & 13.1 & 13.2 \\
     & 500 & 100 & 0.21 & 0.19 & 13.0 & 13.3 \\
     & 500 & 500 & 4.2 & 3.5 & 13.8 & 13.4 \\
     & 2,000 & 20 & 0.61 & 0.57 & 53.7 & 54.8 \\
     & 2,000 & 100 & 0.87 & 0.81 & 45.6 & 46.1 \\
     & 2,000 & 500 & 12.6 & 11.0 & 48.6 & 49.7 \\
     & 8,000 & 20 & 2.8 & 2.7 & 216 & 211 \\
     & 8,000 & 100 & 3.2 & 3.0 & 180 & 182 \\
     & 32,000 & 20 & 11.3 & 10.7 & 988 & 1,015 \\
     & 32,000 & 100 & 11.3 & 10.7 & 655 & 688 \\
    Regression & 500 & 20 & 0.08 & 0.08 & 12.2 & 12.2 \\
     & 500 & 100 & 0.25 & 0.22 & 12.5 & 12.3 \\
     & 500 & 500 & 4.3 & 3.3 & 12.2 & 12.1 \\
     & 2,000 & 20 & 0.38 & 0.36 & 54.9 & 46.3 \\
     & 2,000 & 100 & 1.0 & 0.92 & 42.9 & 43.0 \\
     & 2,000 & 500 & 16.4 & 13.2 & 45.5 & 45.7 \\
     & 8,000 & 20 & 1.7 & 1.6 & 191 & 193 \\
     & 8,000 & 100 & 3.9 & 3.5 & 163 & 163 \\
     & 32,000 & 20 & 8.5 & 8.1 & 788 & 797 \\
     & 32,000 & 100 & 13.2 & 12.1 & 705 & 685 \\
    \bottomrule
  \end{tabular}
  \caption{Fit-only wall-clock seconds for one CIT and one 100-tree CIF in the
  configuration of \cref{app:scaling-curves}, with adaptive stopping (Adaptive)
  and without stopping (No stopping), on Gaussian predictors with a weak planted
  signal in five of them. Each entry is the median of two timed fits, each
  preceded by one untimed identical fit. The cells with $n \ge 8{,}000$ and
  $p=500$ were not run. Tree depth is unbounded, so the forest times are not
  comparable with the depth-3 forests of the head-to-head timing.}
  \label{tab:scaling-curves}
\end{table}
